\originalchapter{离散对称与费米场路径积分}
\label{ch:fermion-path}

连续 Lorentz 变换只连接具有相同空间定向和时间方向的参考系.空间反演、时间反演以及粒子与反粒子的交换,必须另行讨论.它们既帮助判断允许的相互作用,也检查散射振幅的物理解释.另一方面,费米场的反对易统计不能由普通数值函数的路径积分自动产生;需要一种把交换负号写进代数的变量.本章先说明离散对称怎样作用于 Dirac 场,再建立这种 Grassmann 积分.

\originalsection{空间反演与电荷共轭}

记 $x_P=(t,-\mathbf x)$.空间反演 $P$ 是酉算符;在自由 Dirac 理论中可取
\begin{equation}
 P\psi(t,\mathbf x)P^{-1}=\eta_P\gamma^0\psi(t,-\mathbf x),
 \qquad |\eta_P|=1.
 \label{eq:parity-spinor}
\end{equation}
$\eta_P$ 是约定依赖的内禀相位.空间导数在 $x_P$ 处反号,而 $\gamma^0\gamma^i\gamma^0=-\gamma^i$,所以变换后的场满足同一 Dirac 方程.伴随场变为 $\eta_P^*\bar\psi(x_P)\gamma^0$,故标量 $\bar\psi\psi$ 保持不变,矢量流则满足
\begin{equation}
 j^0(t,\mathbf x)\mapsto j^0(x_P),
 \qquad j^i(t,\mathbf x)\mapsto-j^i(x_P).
\end{equation}
对 $\ii\bar\psi\gamma^5\psi$,由 $\gamma^0\gamma^5\gamma^0=-\gamma^5$ 可知它在 $P$ 下变号,因而称为赝标量.$P$ 还交换 $\psi_L$ 与 $\psi_R$;单个 Weyl 理论通常不具有这样的空间反演对称.

电荷共轭 $\mathsf C$ 在场算符 Hilbert 空间中也是酉变换,定义为
\begin{equation}
 \mathsf C\psi(x)\mathsf C^{-1}=\eta_C C\bar\psi^T(x),
 \qquad C=\ii\gamma^2\gamma^0.
\end{equation}
这里右边出现了场的共轭,不意味着 $\mathsf C$ 是反酉算符:它交换独立的 $b,d$ 模,且不把外部复数系数 $\ii$ 变成 $-\ii$.这与下一节的时间反演不同.

对局域双线性,可先分开两个点避免同点接触项,再正规序.电荷共轭后把两个费米场重新排列,会产生一个负号.配合 $C(\gamma^\mu)^TC^{-1}=-\gamma^\mu$,得到
\begin{equation}
 :\bar\psi\psi:\ \xrightarrow{\mathsf C}\ :\bar\psi\psi:,
 \qquad :\bar\psi\gamma^\mu\psi:\ \xrightarrow{\mathsf C}\ -:\bar\psi\gamma^\mu\psi:.
\end{equation}
也可从电荷算符直接检查:$b^\dagger b-d^\dagger d$ 在 $b\leftrightarrow d$ 下变号.下一章的电磁耦合为 $-e j^\mu A_\mu$,因此同时令 $A_\mu\mapsto-A_\mu$,整个耦合才在 $\mathsf C$ 下不变.电子变成正电子时,它所产生的电磁场也反向.

\originalsection{时间反演为何是反酉的}

考虑 $\ii\partial_t\ket{\Psi(t)}=H\ket{\Psi(t)}$.只把 $t$ 换成 $-t$ 而不共轭 $\ii$,方程会得到 $-H$,与正能量谱不相容.因此保持能量的时间反演应是反酉算符 $\mathsf T$,满足
\begin{equation}
 \mathsf T\ii\mathsf T^{-1}=-\ii,
 \qquad\braket{\mathsf T\Phi}{\mathsf T\Psi}=\braket{\Psi}{\Phi}.
\end{equation}
反酉变换对复线性组合共轭系数.对 Dirac 表象中的\emph{数值旋量波函数},可写
\begin{equation}
 \psi_T(t,\mathbf x)=B_T\psi^*(-t,\mathbf x),
 \qquad B_T=\ii\gamma^1\gamma^3.
 \label{eq:time-reversal-wave}
\end{equation}
这条波函数公式与场算符公式 $\mathsf T\psi(t,\mathbf x)\mathsf T^{-1}=B_T\psi(-t,\mathbf x)$ 要区分:后者的复数系数共轭已经由 $\mathsf T$ 的反酉性实施,不能再无条件在右边多取一次共轭.

由矩阵表示可核对 $B_T(\gamma^0)^*B_T^{-1}=\gamma^0$、$B_T(\gamma^i)^*B_T^{-1}=-\gamma^i$.对原 Dirac 方程取复共轭,再以 $t\mapsto-t$ 改变时间导数,便证明式 \eqref{eq:time-reversal-wave} 满足同一方程.动量反向,自旋也反向,因而两者点积即螺旋度不变.$B_TB_T^*=-I_4$,所以单个自旋 $1/2$ 波函数经历两次时间反演得到负号;这正是量子力学中半整数自旋时间反演结构的相对论延续.

矢量流在 $\mathsf T$ 下的时间分量不变、空间分量反号,符合密度不变而速度反向的解释.电磁势应取 $A^0(t,\mathbf x)\mapsto A^0(-t,\mathbf x)$、$\mathbf A(t,\mathbf x)\mapsto-\mathbf A(-t,\mathbf x)$,从而电场在 $T$ 下不变、磁场变号.结合空间反演,电场在 $P$ 下变号、磁场不变;在 $\mathsf C$ 下两者都变号.这些可直接由 $\mathbf E=-\nabla A^0-\dot{\mathbf A}$ 和 $\mathbf B=\nabla\times\mathbf A$ 检查.

\begin{remark}
本章对自由 Dirac 理论和矢量电磁耦合逐项检查了 $C,P,T$.一般局域、Lorentz 协变、正能量并具有适当 Hilbert 空间结构的量子场论满足 CPT 定理,即三者的联合是对称;其完整证明需要解析延拓和局域性工具,超出这里的推导范围.CPT 并不意味着三种单独对称都成立,弱相互作用就是不能由本章的 QED 检查推广的情形.
\end{remark}

\originalsection{Grassmann 变量及其积分}

普通积分变量彼此交换,无法在交换两个费米外腿时自动给出负号.设 $\theta_1,\ldots,\theta_n$ 是形式生成元,满足
\begin{equation}
 \theta_i\theta_j=-\theta_j\theta_i,
 \qquad \theta_i^2=0.
\end{equation}
由这些生成元的有限线性组合构成 Grassmann 代数.一个变量的函数只有 $f(\theta)=a+b\theta$;多个变量的函数是次数不超过 $n$ 的多项式.奇次项与奇次项交换时变号,偶次项与所有项交换.路径积分中的 Grassmann 场是取值于这种代数的形式场,并不是逐点可观察的随机实数.

\begin{definition}[Berezin 积分]
规定 $\int\dd\theta\,1=0$、$\int\dd\theta\,\theta=1$,并对系数线性延拓.采用左导数,$\partial_{\theta_i}\theta_j=\delta_{ij}$,且奇变量乘积求导满足
\begin{equation}
 \partial_{\theta_i}(AB)=(\partial_{\theta_i}A)B+(-1)^{|A|}A\partial_{\theta_i}B,
\end{equation}
其中 $|A|=0,1$ 分别表示偶、奇次数.
\end{definition}

因此积分与取一次 Grassmann 导数等价,却与普通面积积分含义不同.平移 $\theta\mapsto\theta+\epsilon$ 不改变最高次系数,所以积分平移不变.若 $\theta=a\eta$,要求 $\int\dd\theta\,\theta=1$ 给出 $\dd\theta=a^{-1}\dd\eta$,这是普通变量 Jacobian 的倒数.对 $n$ 个变量推广为 $\dd^n\theta=(\det M)^{-1}\dd^n\eta$.这一反常的 Jacobian 正是费米 Gaussian 积分出现行列式而非其倒数的原因.

$\bar\theta_i$ 是独立 Grassmann 生成元;积分时不对它施加复共轭约束.固定测度方向,使单变量 $\int\dd\bar\theta\dd\theta\,\theta\bar\theta=1$.于是
\begin{equation}
 \int\dd\bar\theta\dd\theta\,\e^{-\bar\theta a\theta}
 =\int\dd\bar\theta\dd\theta(1+a\theta\bar\theta)=a.
\end{equation}
交换测度顺序会同时改变整体符号,故必须始终保持同一约定.多变量取配对有序测度 $\dd\bar\theta_n\dd\theta_n\cdots\dd\bar\theta_1\dd\theta_1$,则
\begin{equation}
 \int D\bar\theta D\theta\,\e^{-\bar\theta_i A_{ij}\theta_j}=\det A.
 \label{eq:grassmann-gaussian}
\end{equation}
指数中只有选中每一个 $\theta_j,\bar\theta_i$ 各一次的项能在积分后保留.这些项由排列 $\pi$ 标记,重排变量产生 $\operatorname{sgn}(\pi)$,系数是 $\prod_i A_{i\pi(i)}$;总和恰是行列式.这个有限维推导说明了费米统计从何处进入积分,并不依赖给积分一个概率分布解释.

\originalsection{外源、传播子与 Wick 负号}

加入独立奇源 $\eta,\bar\eta$.对可逆矩阵 $A$,平移变量完成平方,有
\begin{align}
 &-\bar\theta A\theta+\bar\eta\theta+\bar\theta\eta
 =-(\bar\theta-\bar\eta A^{-1})A(\theta-A^{-1}\eta)
   +\bar\eta A^{-1}\eta,\\
 &\int D\bar\theta D\theta\,
 \e^{-\bar\theta A\theta+\bar\eta\theta+\bar\theta\eta}
 =\det A\,\e^{\bar\eta A^{-1}\eta}.
\end{align}
源也必须是奇变量,才能让整个指数为偶数,并使不同源微分的交换负号与算符统计一致.

把有限维指标扩充为旋量指标和时空点,在时空离散化的意义下定义自由 Dirac 生成泛函
\begin{equation}
 Z[\bar\eta,\eta]=\frac{\int D\bar\psi D\psi\,
 \exp\{\ii\int\dd^4x[\bar\psi K\psi+\bar\eta\psi+\bar\psi\eta]\}}
 {\int D\bar\psi D\psi\,\exp\{\ii\int\dd^4x\bar\psi K\psi\}},
 \qquad K=\ii\slashed\partial-m.
\end{equation}
逆算符取真空 Feynman 边界条件,记为 $K_F^{-1}$.平移 $\psi\mapsto\psi+K_F^{-1}\eta$、$\bar\psi\mapsto\bar\psi+\bar\eta K_F^{-1}$ 后得到
\begin{equation}
 Z[\bar\eta,\eta]
 =\exp\!\left[-\ii\int\dd^4x\dd^4y\,
       \bar\eta(x)K_F^{-1}(x,y)\eta(y)\right]
 =\exp\!\left[-\int\dd^4x\dd^4y\,\bar\eta(x)S_F(x-y)\eta(y)\right].
\end{equation}
最后一步用 $S_F=\ii K_F^{-1}$.Grassmann 源的微分次序必须与插入算符次序一起规定;不能像玻色源那样任意交换.采用已定义的左导数,先对 $\bar\eta_a(x)$ 再对 $\eta_b(y)$ 求导,在零源处有
\begin{equation}
 \left.\frac{\delta^L}{\delta\eta_b(y)}\frac{\delta^L}{\delta\bar\eta_a(x)}Z\right|_{\eta=\bar\eta=0}=-S_{F,ab}(x-y).
\end{equation}
因此提取排列为 $\psi_a(x)\bar\psi_b(y)$ 的两点插入时,再乘 $1/\ii^2=-1$,恢复 $S_F$.这个负号可从指数 $-\bar\eta S_F\eta$ 直接核对;更多插入同样须按固定源导数次序保留交换号.比如展开四点函数,先保留原排列 $\psi_1\bar\psi_2\psi_3\bar\psi_4$,可得
\begin{equation}
 \bra0T\psi_1\bar\psi_2\psi_3\bar\psi_4\ket0
 =S_{12}S_{34}-S_{14}S_{32}.
 \label{eq:fermion-wick-four}
\end{equation}
每个 $S_{ij}$ 带相应旋量指标.第二项为使配对算符相邻而产生了奇数次费米交换.这个具体例子把 Grassmann 源展开与前章算符时间排序连接起来.

\originalsection{闭合费米圈的负号}

若背景场使 $K$ 变成 $K+V$,积分掉 Dirac 场得到 $\det(K+V)$.先在有限维正规化中使用
\begin{equation}
 \det(K+V)=\det K\,
 \exp\left[\Tr\ln(1+K^{-1}V)\right],
 \quad
 \Tr\ln(1+X)=\sum_{n\geq1}\frac{(-1)^{n+1}}n\Tr X^n.
\end{equation}
$\Tr$ 同时包含时空积分与旋量迹,$\tr$ 只表示旋量迹.每项 $\Tr(K^{-1}V)^n$ 是一条经过 $n$ 个顶点并闭合的有向费米线.普通复玻色 Gaussian 给出 $(\det K)^{-1}$,费米 Gaussian 给出 $\det K$;取对数后两者符号相反.这与 Wick 展开中把闭合费米链首尾接合所需的一次额外交换相同.因此 Feynman 规则规定每个闭合费米圈额外乘 $-1$.

这里的行列式和迹在连续场论中仍需正规化;尚未声称它们是收敛的无穷维普通行列式.下一章先在树图中应用这些规则,第十二章再说明环积分为什么需要正规化,以及这些形式表达式怎样与测得的参数联系.
